\chapter{Tabla de pines}\label{anx:pines}

\section{Pines de datos de las caras}

\begin{tabularx}{\linewidth}{@{}>{\raggedright\arraybackslash}p{0.3\linewidth} X X X@{}}
\toprule
\textbf{Señal} & \textbf{Esquema PCB\_V2 (netlist)} & \textbf{Poste funcionando} & \textbf{Driver en WLED} \\ \midrule
Cara A (\texttt{GPIO\_A} → R1 → \texttt{D1} → A1) & GPIO\dato{pinesEsquema.GPIO_A} & GPIO\dato{pinesPosteActual.caraA} & RMT (\texttt{drv: 0}) \\
Cara B (\texttt{GPIO\_B} → R2 → \texttt{D2} → B1) & GPIO\dato{pinesEsquema.GPIO_B} & GPIO\dato{pinesPosteActual.caraB} & I2S (\texttt{drv: 1}) \\
Reserva (\texttt{GPIO\_C} → B2) & GPIO\dato{pinesEsquema.GPIO_C} & --- & --- \\
\bottomrule
\end{tabularx}

En WLED, las dos salidas empiezan en el LED 0 sea cual sea el pin: la cara B repite a la cara A (sección~\ref{sec:salidas-wled}). Los pines del poste que funciona están confirmados: la ESP32 va en una placa de borneras aparte y sus GPIO\dato{pinesPosteActual.caraA} y \dato{pinesPosteActual.caraB} llegan por cable a la bornera B2 (sección~\ref{sec:pines}). Con la ESP32 soldada en la \dato{pcbVersion}, los datos salen por GPIO\dato{pinesEsquema.GPIO_A} y GPIO\dato{pinesEsquema.GPIO_B}.


\section{Alimentación y otros pines de la ESP32}

\begin{tabularx}{\linewidth}{@{}l l X@{}}
\toprule
\textbf{Pin del módulo} & \textbf{Red} & \textbf{Uso} \\ \midrule
21 (5V) & \texttt{+5V\_ESP} & Alimentación desde el LM2596 (\dato{reguladorSalida}). \\
22, 23, 24, 44 (GND) & GND & Tierra común con la fuente y las caras. \\
GPIO0 (botón BOOT) & --- & Modo de grabación. En WLED es el botón 0: de 6 a 9~s abre la red de configuración; más de 10~s borra la memoria de archivos. \\
\bottomrule
\end{tabularx}

\nota{Si puedes elegir, no uses como datos los GPIO 19 y 20 (USB), 26 a 37 (memoria flash y PSRAM de la N16R8) ni los pines de arranque 0, 3, 45 y 46. GPIO3 funciona en el poste actual como datos de la cara B; como es un pin de arranque, no lo fuerces a un nivel fijo con resistencias externas.}
